\chapter{Datasets}

We run our experiments on two datasets: MUSIN-G\cite{musing} and what we call BCMI-training\cite{bcmi}.
Both contain joint EEG and music recordings, but while MUSIN-G contains natural music, BCMI-training contains computer-generated music.
We choose the two datasets to test our hypothesis both on generated but uniform and natural music.

\section{MUSIN-G}

It is a dataset of joint music and EEG recordings of 20 non-musician subjects, aged 22-28 years\cite{musing}.
Data was acquires using high-density Hydrocel Geodesic Sensor Net.
Each subject was presented with 12 songs of various genres, ranging from Indian Folk to Electronic Dance.

\begin{table}[h]
\centering
\begin{tabular}{|l|l|l|c|}
\hline
Song Name & Artist & Genre & Length (s) \\
\hline
Trip to the Lonely Planet & Mark Alow & Deep House & 125 \\
Sail & AWOLNATION & Electric Pop & 114 \\
Concept 15 & Kodomo & Electronic & 132 \\
Aurore & Claire David & New Age & 111 \\
Proof & Idiotape & Electronic & 124 \\
Glider & Tycho & Slow Electronic & 100 \\
Raag Bihag & Pandit Hariprasad Chaurasia & Hindustani Classical & 116 \\
Albela Sajan & Shankar Mahadevan & Indian Semi-Classical & 121 \\
Mor Bani Thanghat Kare & Aditi Paul & Indian Folk & 126 \\
Fallin & Dr. SaxLove & Soft Jazz & 129 \\
Master of Running & Rickeyabo & Goth Rock & 113 \\
JB & Nobody.one & Progressive Instrumental & 117 \\
\hline
\end{tabular}
\caption{MUSIN-G dataset songs}
\label{tab:musing-songs}
\end{table}



\section{BCMI-training}

BCMI-training is a dataset of joint music and EEG recordings of 10 participants aged 19-30.
It is part of a series of datasets collected during studies investigating human reactions to affective musical stimuli and whether reported emotionality can be influenced by music\cite{nicolaou}.
The music played in this study is fully computer-generated, and no track repeats across the dataset.
The musical snippets are tagged with emotionality labels (9 possibilities on valence/arousal grid) which are inputs to the music generation algorithm.
In the study, subjects listened to a musical snippet of about 20 seconds in length and immediately afterward to another snippet of different emotionality.
For our purposes, we split such trials into two separate recordings, one for each snippet.

\section{Dataset processing}

The raw datasets are preprocessed for the use in our experiments.
To efficiently construct various variants of the datasets, including different preprocessing methods, we created a library for dataset handling and processing.
The library itself uses libraries like "librosa" and "essentia" for audio processing and "mne" for EEG processing.
The library allows for loading of the two datasets to a common format, lazy evaluation of preprocessing steps, defines a standard on disk storage format and simplifies dataset operations like filtering. 
It also allows to combine datasets, but we didn't end up using this feature.
We show below how a dataset is loaded for a study using our library.

We unify the various analyzed datasets under the `EEGMusicDataset` class, which provides a common interface to access the data and metadata of the recordings.
Then, various preprocessing functions are applied to the data; their evaluation is lazy.
Lastly, the `StratifiedSamplingDataset` wrapper is applied, which samples the data in a stratified way, as described in the methods chapter.

\begin{verbatim}
ds = EEGMusicDataset.load_ondisk(Path("..."))

ds = MappedDataset(ds, lambda x: 
  ica_band_power_trial(
    rereference_trial(
      filter_trial(
        notch_filter_trial(x, 50.0), l_freq=0.5, h_freq=50.0
      )
    ),
    n_components=12,
    bands = [(0.5, 4), (4, 8), (8, 13), (13, 30), (30, 45)],
    window_sec= 2.0,
    hop_sec = 0.1,
  )
)

ds = StratifiedSamplingDataset(ds, n_strata=10, trial_length_secs=Fraction(3, 1))

\end{verbatim}

% musing


% grupa bcmi:
%  - opis eksperymentu
%  - używamy treningowego, choć można połączyć